\documentclass{article}
\usepackage{circuitikz}
\begin{document}

\begin{figure}[ht]
\centering
\begin{circuitikz}
  \draw (0,0) node[pmos, anchor=D] (p1) {};
  \draw (0,0) node[nmos, anchor=D] (n1) {};
  \draw (p1.D) -- (n1.D);
  \draw (p1.S) node[vdd] {$V_{DD}$};
  \draw (n1.S) node[ground] {};
  \draw (p1.G) -- (n1.G);
  \draw ($(p1.G)!0.5!(n1.G)$) to[short, *-o] ++(-1,0) node[left] {$V_{in}$};
  \draw (0,0) to[short, *-o] (1.5,0) node[right] {$V_{out}$};
\end{circuitikz}
\caption{CMOS inverter built from a PMOS pull-up and an NMOS pull-down transistor.}
\label{fig:inverter}
\end{figure}

\end{document}
